\section{The Confusion of Levels}

In \emph{The Reign of Quantity}, \textbf{Guenon} has a chapter on the consequences of the failure of modern man to distinguish between the psychic and the spiritual without, however, clearly delineating the distinction. Man is a tripartite being, consisting of body, soul, and spirit. Each level has its own proper functions, which mimic the corresponding function of other levels, thus making possible confusion between the levels. Chart \ref schematizes the functions.

\begin{table}[h]
\begin{tabular}{r*{3}{p{.25\textwidth}}}\toprule
Levels & \multicolumn{3}{l@{}}{Functions} \\\midrule
~ & \textbf{Feeling} & \textbf{Knowing} & \textbf{Willing}\\\cmidrule{2-4}
Body (\textbf{Hylic}) & Attraction/Aversion & Instinct & Movement \\
Soul (\textbf{Psychic}) & Emotion & Thoughts & Desire \\
Spirit (\textbf{Pneumatic}) & Communion & Intuition & Conscience \\\bottomrule
\end{tabular}
\end{table}

In the \textbf{True Man}, the psychic level mediates the hylic and pneumatic levels. However, the undeveloped man, common to the modern world, identifies fully with the psychic level (or even worse, the hylic level). He regards his emotions, thoughts, and desires as ``his own'', unaware of their true source and significance. Since every emotion, thought, and desire engenders its opposite, the psychic man's identity is inconsistent, varying with their changes.

For the pneumatic man, the functions have no opposites. Unlike the psychical emotions which have their negative counterparts, the higher emotions do not. There is a direct contact, or communion with higher states. Intuition provides direct and unmediated understanding (gnosis) of higher reality. Willing is not plagued by doubt, but the pneumatic man has \textbf{True Will} and knows what to do. He knows the virtues directly -- prudence, courage, justice, temperance --- and understands them as power. His conscience does not experience moral failure as guilt against some moral standard, but rather as shame at failing to live up to his own. The development of these functions are forms of Bhakti, Jnana, and Karma Yoga.

The danger comes when ideas appropriate to the pneumatic man are adopted by the psychic man. For example, Rousseau writes:

\begin{quotex}
What God wills a man to do, he does not tell him through another man; He tells it to him Himself, and writes it in the bottom of his heart.
\end{quotex}


This opposes a leveling democratism to legitimate hierarchy, which is freely accepted. It engenders disloyalty, since every man is his own highest authority. The symbol of the ``Heart'', which should represent the spiritual functions, is instead debased into little more than a sentimental meaning.

The number of such examples can be multiplied indefinitely. So, instead, we will generalize how the confusion of levels leads to the three fundamental philosophical deformations in the modern world. They are characterized by being anti-Tradition. Although they are logically inconsistent and incompatible with each other, in practice they unite against Tradition since none of them bases itself on higher principles; thus their opposition to Principle in held in common.

\begin{itemize}
\item \textbf{Sentimentalism} is the result of confusion of the feeling function
\item \textbf{Rationalism} is the result of the confusion of the thinking function
\item \textbf{Moralism} is the result of the confusion of the willing function
\end{itemize}

Reference: \emph{The Great Triad} by Rene Guenon. The chapters on Body, Mind, Soul and on True Man


\flrightit{Posted on 2010-09-20 by Cologero}

\begin{center}* * *\end{center}

\begin{footnotesize}
\begin{sffamily}
\texttt{Charles Upton on 2010-09-20 at 00:54 said:}

Correct.

\hfill

\texttt{Will on 2010-09-20 at 22:06 said:}

Rama Coomaraswamy offers some good insight on this matter here:
\url{http://www.sacredweb.com/online\_articles/sw9\_coomaraswamy.html}



\hfill

\end{sffamily}
\end{footnotesize}

